% Part: normal-modal-logic
% Chapter: axioms-systems
% Section: provability-properties

\documentclass[../../../include/open-logic-section]{subfiles}

\begin{document}

\olfileid{nml}{prf}{prp}

\olsection{\usetoken{S}{derivability} లక్షణాలు}

\begin{prop}\ollabel{prop:derivabilityfacts}
  $\Sigma$ ఒక మోడల్ వ్యవస్థ, $\Gamma$ మోడల్
  \tetoken{సూత్రాల}{!!{formula}s} సమితి అనుకుందాం. అప్పుడు కింది లక్షణాలు వర్తిస్తాయి:
  \begin{enumerate}
  \item\ollabel{prop:derivabilityfacts-monotonicity} \emph{ఏకదిశత}:
    $\Gamma \Proves[\Sigma] !A$ మరియు $\Gamma \subseteq \Delta$ అయితే,
    $\Delta \Proves[\Sigma] !A$;
  \item\ollabel{prop:derivabilityfacts-reflexivity}
    \emph{స్వావర్తనత్వం}: $!A \in \Gamma$ అయితే, $\Gamma
    \Proves[\Sigma] !A$;
  \item\ollabel{prop:derivabilityfacts-cut} \emph{కట్}: $\Gamma
    \Proves[\Sigma] !A$ మరియు $\Delta \cup \{!A\} \Proves[\Sigma] !B$
    అయితే, $\Gamma \cup \Delta \Proves[\Sigma] !B$;
  \item\ollabel{prop:derivabilityfacts-deduction} \emph{నిగమన సిద్ధాంతం}:
    $\Gamma \cup \{!B\} \Proves[\Sigma] !A$ అయితే, అప్పుడు మాత్రమే
    $\Gamma \Proves[\Sigma] !B \lif !A$;
  \item \ollabel{prop:derivabilityfacts-ruleT}% \emph{Rule T}: If
    $\Gamma \Proves[\Sigma] !A_1$ మరియు \dots{} మరియు $\Gamma
    \Proves[\Sigma] !A_n$ అయ్యి, $!A_1 \to (!A_2 \lif \cdots (!A_n \lif
    !B)\cdots)$ సర్వసత్య నిదర్శనమైతే, $\Gamma
    \Proves[\Sigma] !B$.
  \end{enumerate}
\end{prop}

నిరూపణ సులభమైన వ్యాయామం. \olref{prop:derivabilityfacts-ruleT},
అంటే \olref{prop:derivabilityfacts}లోని ఆ భాగం వల్ల, ఉదాహరణకు,
$\Gamma \Proves[\Sigma] !A \lor !B$ మరియు
$\Gamma \Proves[\Sigma] \lnot !A$ అయితే,
$\Gamma \Proves[\Sigma] !B$ అని వస్తుంది. అలాగే, ఇకపై
$\Gamma \cup \{ !A\} \Proves[\Sigma] !B$ బదులు
$\Gamma, !A \Proves[\Sigma] !B$ అని రాస్తాం.

\begin{defn}
  $\Gamma$ సమితి, $\Sigma$ వ్యవస్థకు సాపేక్షంగా
  \emph{నిగమన పరంగా సంవృతం} అంటే, $\Gamma \Proves[\Sigma] !A$
  అయినప్పుడల్లా $!A \in \Gamma$ కావడం.
\end{defn}

\end{document}
